\documentclass[UTF8,fontset=none,11pt,a4paper]{ctexart}
\usepackage{ipara-handbook}
\handbooksetup{NET-B06 MTU × Segmentation}{408 · Network Internal Bridge}{Message · MSS · Packet · Frame}
\begin{document}
\handbooktitle{MTU × Segmentation}{路径尺寸怎样约束传输层分段}{408 · NET-B06}

\begin{corebox}{Mother Interface}
\[
\boxed{Application\ bytes\xrightarrow{TCP\ MSS}segments
\xrightarrow{IP\ MTU}packets\xrightarrow{link\ framing}frames}
\]
Segmentation 主动把 byte stream 切成传输单元；fragmentation 被动把过大的 IP packet 适配链路 MTU。两者的 Owner、首部、序号与重组位置不同。
\end{corebox}

\begin{methodbox}{Owners 与停止边界}
NET06 Owns TCP stream/segment 与 MSS；NET04 Owns IP packet、MTU、fragmentation 和 PMTU。本桥 Owns 尺寸预算与失败路由。Ethernet/PPP frame 开销回到 NET02，具体可靠重传回到 NET-B03。
\end{methodbox}

\section{四层边界不能混算}

\begin{center}\small
\begin{tabularx}{\linewidth}{L{2.2cm}Y Y Y}
\toprule
对象 & 尺寸参数 & 谁切/封装 & 谁重组/解释\\
\midrule
Application message & 应用自定义 & application & application\\
TCP segment & MSS 限制 TCP data & source TCP & destination TCP 按 byte SEQ 重排\\
IP packet/fragment & MTU 限制整个 IP packet & source；IPv4 router 在允许时也可分片 & final IP destination\\
Link frame & frame payload/format & each-hop link adapter & next-hop link endpoint\\
\bottomrule
\end{tabularx}
\end{center}

TCP 是 byte stream，不保留 application write boundary；一个 write 可成为多个 segments，多个 write 也可能合并。MSS 是每个 TCP segment 的最大 data 部分，不含 TCP/IP headers。经典无选项 IPv4 预算常写：
\[
MSS\le PMTU-H_{IP}-H_{TCP}.
\]
若 PMTU 为 1500 B、IPv4/TCP 首部各 20 B，则 MSS 为 1460 B；有 IP/TCP options、IPv6 extension headers 或隧道外层时必须重新扣除实际开销。

\section{为什么优先分段而不是依赖 IP 分片}

TCP 在源端按合适 MSS 生成不超过 PMTU 的 packet，可避免中途 IP fragmentation。若一个大 IPv4 packet 已形成且输出 MTU 更小，允许时 router 可分片；所有 fragments 只在最终 IP destination 重组，router 不逐段恢复 TCP。任一 fragment 丢失会使该原 packet 无法交给 TCP，TCP 最终重传的通常是相关字节数据，可能再次生成新的 packet，而不是让中间 router 重传一个 fragment。

IPv6 router 不分片：收到 Packet Too Big 后，source 调整 packet size/PMTU，并在必要时由 source 使用 Fragment header。故“IPv6 不分片”缺少 Owner，是错误压缩。

\section{封装开销会改变有效 MTU}

隧道在原 packet 外再增加 outer header。若物理 egress MTU 不变，inner packet 可用预算随外层开销减少：
\[
MTU_{inner}\le MTU_{egress}-H_{outer}.
\]
Mobile IP、VPN 或其他隧道题若只沿用原 MSS，可能在封装后超限。链路 frame 的 header/trailer 通常不计入 IP MTU，但题设若给的是“最大帧长”而不是 MTU，必须先扣链路开销得到网络层 payload budget。

\section{数值母例}

应用写入 4000 B，TCP MSS=1460 B，则可形成 data 长度 $1460,1460,1080$ B 的三个 TCP segments，各自拥有 TCP header 与连续 byte SEQ。若其中一个 1500 B IPv4 packet 之后遇到 MTU 620 B，IPv4 header 20 B 且允许 router 分片，则非末片 payload 取不超过 600 B 的 8-byte 整数倍；这是对整个 IP payload（含 TCP header/data）的 IP fragmentation，不是 TCP 再按 MSS 分段。

\section{调用协议与 First Divergence}

\begin{enumerate}
\item 写题设给的是 application bytes、MSS、IP MTU 还是 maximum frame size；
\item 逐层扣除对应 header，绝不跨层重复扣或漏扣；
\item 判断是 source TCP segmentation、IPv4 fragmentation 还是 IPv6 PMTU 分支；
\item 分段按 TCP byte SEQ 推进；分片按 Identification/Offset/MF 推进；
\item 明确重组 Owner，并检查隧道外层是否再次消耗 MTU。
\end{enumerate}

\begin{warnbox}{Anti-Bridge}
把 MSS 说成最大 TCP segment 总长度，是漏了首部；把 fragment offset 当 TCP SEQ，是混淆两个坐标系；router 收齐 fragments 再转发，是写错重组 Owner；链路最大帧长直接当 IP MTU，是忽略 frame overhead；IPv6 router 主动分片，是写错责任边界。
\end{warnbox}

\begin{corebox}{Compression}
当前对象是哪一层？尺寸上限包含哪些首部？谁切分、用什么坐标、谁重组？外层封装是否缩小 inner budget？
\end{corebox}
\end{document}
